\documentclass[10pt,a4paper]{article}
\usepackage[margin=0.85in]{geometry}
\usepackage[T1]{fontenc}
\usepackage[utf8]{inputenc}
\usepackage{lmodern,amsmath,amssymb,graphicx,enumitem,microtype,fancyhdr,xcolor,needspace}
\usepackage[hidelinks]{hyperref}
\definecolor{lightgrey}{gray}{0.94}
\definecolor{midgrey}{gray}{0.35}
\setlength{\parindent}{0pt}
\setlength{\parskip}{5pt}
\setlist[enumerate]{itemsep=5pt,topsep=4pt}
\pagestyle{fancy}\fancyhf{}\renewcommand{\headrulewidth}{0pt}
\fancyfoot[L]{\footnotesize\color{midgrey}Arij Asad}
\fancyfoot[C]{\footnotesize\color{midgrey}Edexcel companion practice}
\fancyfoot[R]{\footnotesize\color{midgrey}\thepage}
\newcounter{question}
\newenvironment{question}[1]{\refstepcounter{question}\Needspace{5\baselineskip}\par\medskip\textbf{\thequestion.\ #1}\par}{\par\vspace{12mm}}
\begin{document}
\hypersetup{pdftitle={Statistics and Mechanics Year 2 - Original practice and solutions},pdfauthor={Arij Asad}}
\begin{minipage}[c]{0.78\textwidth}{\Large\bfseries Statistics and Mechanics Year 2}\par Original practice check\end{minipage}\hfill\begin{minipage}[c]{0.17\textwidth}\IfFileExists{PS_logo.png}{\includegraphics[width=\linewidth]{PS_logo.png}}{\rule{0pt}{16mm}}\end{minipage}

\medskip\fcolorbox{black!20}{lightgrey}{\parbox{\dimexpr\linewidth-2\fboxsep-2\fboxrule\relax}{Three short checks on the normal distribution, projectiles and moments. A calculator may be used. Take gravitational acceleration as 9.8 metres per second squared. Attempt the questions before reading the solutions. These original questions sample selected skills; they are not an official Edexcel paper.}}

\begin{question}{Normal distribution}
\(X\) is normally distributed with mean \(50\) and standard deviation \(4\). Find \(P(X>54)\), giving your answer to three significant figures.
\end{question}
\begin{question}{Projectiles}
A particle is projected from level ground at \(20\,\mathrm{m\,s^{-1}}\), at \(30^\circ\) above the horizontal. Neglect air resistance and take \(g=9.8\,\mathrm{m\,s^{-2}}\). Find its time of flight and horizontal range when it returns to the same level.
\end{question}
\begin{question}{Moments and equilibrium}
A uniform horizontal beam \(AB\) is \(4\,\mathrm{m}\) long and weighs \(40\,\mathrm{N}\). It is supported at \(A\) and \(B\). A further downward load of \(20\,\mathrm{N}\) acts \(1\,\mathrm{m}\) from \(A\). Find the upward reactions at the two supports.
\end{question}
\Needspace{7\baselineskip}\textbf{Find the mistake}\par
Events \(A\) and \(B\) satisfy \(P(A)=0.4\), \(P(B)=0.5\) and \(P(A\cap B)=0.3\). A student writes \(P(A\mid B)=P(A)=0.4\). Explain why this is wrong and find the correct conditional probability.

\vfill\small Further practice: \href{https://maths.arijasad.com/tmua-paper-archive.html}{TMUA papers and solutions} and \href{https://maths.arijasad.com/maths-topic-practice.html}{maths topic guides}.
\newpage
{\Large\bfseries Hints and worked solutions}\par\medskip
\Needspace{8\baselineskip}\subsection*{1. Normal distribution}
\textit{Hint: Standardise using the standard deviation, then choose the upper tail.}\par
With \(Z=(X-50)/4\), the threshold \(X=54\) corresponds to \(Z=1\).

Hence \(P(X>54)=P(Z>1)=1-\Phi(1)\approx0.158655\), where \(\Phi\) is the standard normal cumulative distribution function.

\textbf{Answer:} \(1-\Phi(1)\approx0.159\) (three significant figures).

\Needspace{8\baselineskip}\subsection*{2. Projectiles}
\textit{Hint: Resolve the initial velocity into components. Use zero vertical displacement to find the non-zero flight time.}\par
The initial horizontal and vertical velocity components are \(20\cos30^\circ=10\sqrt3\,\mathrm{m\,s^{-1}}\) and \(20\sin30^\circ=10\,\mathrm{m\,s^{-1}}\).

Taking upward as positive, vertical displacement is \(y=10t-4.9t^2\). At landing, \(y=0\); the non-zero solution is \(t=10/4.9=100/49\,\mathrm{s}\approx2.04\,\mathrm{s}\).

Horizontal velocity is constant, so the range is \((10\sqrt3)(100/49)=1000\sqrt3/49\,\mathrm{m}\approx35.3\,\mathrm{m}\).

\textbf{Answer:} Time \(\dfrac{100}{49}\,\mathrm{s}\); range \(\dfrac{1000\sqrt3}{49}\,\mathrm{m}\) (approximately \(2.04\,\mathrm{s}\) and \(35.3\,\mathrm{m}\)).

\Needspace{8\baselineskip}\subsection*{3. Moments and equilibrium}
\textit{Hint: The beam's weight acts at its midpoint. Taking moments about one support eliminates its reaction.}\par
Let the upward reactions be \(R_A\) and \(R_B\). Vertical equilibrium gives \(R_A+R_B=40+20=60\,\mathrm{N}\).

Taking moments about \(A\), \(4R_B=40(2)+20(1)=100\). Therefore \(R_B=25\,\mathrm{N}\).

It follows that \(R_A=60-25=35\,\mathrm{N}\).

\textbf{Answer:} \(R_A=35\,\mathrm{N}\), \(R_B=25\,\mathrm{N}\).

\Needspace{9\baselineskip}\subsection*{Find the mistake: correction}
\(P(A\mid B)=P(A\cap B)/P(B)=0.3/0.5=0.6\).

The student's equality would hold for independent events. Here \(P(A)P(B)=0.4(0.5)=0.2\), which differs from \(P(A\cap B)=0.3\); the events are not independent.

\Needspace{6\baselineskip}\subsection*{What to practise next}\begin{enumerate}
\item For normal probabilities, draw and shade the required tail before entering values; distinguish variance from standard deviation.
\item For projectiles, keep separate horizontal and vertical equations linked by the same time.
\item For equilibrium, draw every force at its point of application and choose a moment centre that removes an unknown.
\end{enumerate}\end{document}
